% Irodov Part V §5.1 26--35. User PDF pp.199--201, printed 190--192,
% geometric figures 5.4--5.8 and source answers pp.325--326.
\opticsproblem{26}{190--191}
用作图法完成原书图5.4、5.5两组任务：
(a) 已知凹面和凸面球面镜、主轴及焦点，续作指定入射光线；
(b) 已知一对物像共轭点 P、P' 及主轴，确定镜面位置和焦点。
\exsolution
(a) 对球面镜使用两条基本射线：
与主轴平行的近轴入射光线反射后经过焦点（凸镜则其反向延长线经过虚焦点）；
指向曲率中心的光线沿原路返回。
从图5.4中已知焦点 $F$ 和镜顶点 $O$，曲率中心 $C$ 由
$\overline{OC}=2\overline{OF}$ 构造。
先作入射光线与镜面交点处的球面法线 $MC$，再将入射方向绕法线镜像，即得反射光线。
(b) 已知 P 与 P'，光源经镜面顶点反射的傍轴光线满足入射和反射角相等，可由共轭点的角平分线确定镜顶点；
再依有向球面镜公式
\[
\boxed{\frac1s+\frac1{s'}=\frac1f=\frac2R}
\]
确定焦距及曲率中心；对实、虚像用有向距离符号。
\excheck 原题为\emph{图解作图}，上述规则给出可执行的作图步骤，不捏造图5.4或5.5里未标出的数值。

\opticsproblem{27}{191}
凹面镜：(a) 物与像的距离 $l=15$cm、横向放大率
$\beta=-2.0$；
(b) 物体沿主轴移动 $l=5.0$cm 后，放大率由
$\beta_1=-0.50$ 变为 $\beta_2=-0.25$。
分别求镜焦距。
\exsolution
在实物实像约定下
$\beta=-s'/s$，
镜公式 $1/f=1/s+1/s'$。
(a) $s'=2s$，物像轴向间距
$|s'-s|=s=l$，
所以 $s=15$cm、$s'=30$cm，
\[
\boxed{f=\frac{ss'}{s+s'}=10\,\mathrm{cm}}.
\]
(b) 利用
$\beta=-f/(s-f)$，
可解物距
$s=f(1-1/\beta)$。
故两位置
$s_1=3f$、$s_2=5f$，
相距 $2f=l=5$cm，
\[
\boxed{f=2.5\,\mathrm{cm}}.
\]
\excheck 对实像放大率负号不可忽略；镜焦距为正的凹面镜光学绝对值。

\opticsproblem{28}{191}
点光源光强 $I_0=100$cd，距凹面镜顶点
$s=20.0$cm，焦距 $f=25.0$cm，镜面反射率
$\rho=0.80$。
求反射光束等效点像的光强。
\exsolution
球面镜成像公式
$1/s+1/s'=1/f$，
所以
$s'=fs/(s-f)
=(25)(20)/(20-25)=-100$cm，
形成放大率模
$|\beta|=|s'/s|=5$ 的虚像。
镜面反射后的同一束光功率是入射功率的 $\rho$ 倍，
从虚像发出的相同锥形方向固角缩小为原来的
$(s/|s'|)^2$ 倍（相当于等效辐射强度增加 $\beta^2$），所以
\[
\boxed{I'=\rho I_0
\left(\frac{f}{f-s}\right)^2
=0.80(100)(5)^2
=2.0\times10^3\,\mathrm{cd}}.
\]
\excheck 这里 $I'$ 为镜面反射束的\emph{等效像光强}，不是镜子新产生了额外总光通量。

\opticsproblem{29}{191}
从费马原理出发，推导折射率分别为
$n,n'$ 的两种介质之间的球面折射界面，对近轴光线满足的成像公式（界面曲率半径 $R$）。
\exsolution
将界面顶点置于原点，曲率球心在有向 $+x$ 方向为正，近轴表面高度 $h$ 处的轴向弓高
$x(h)\simeq h^2/(2R)$。
设物点、像点的有向轴距分别 $s,s'$，
按费马原理同一物像共轭光束的光程不随交点高度二阶改变。
从物点到界面、从界面到像点的光程展开至 $h^2$：
\[
n\sqrt{(x-s)^2+h^2}
+n'\sqrt{(s'-x)^2+h^2}
=\text{常数}+O(h^4).
\]
令 $h^2$ 系数为零，整理得到
\[
\boxed{\frac{n'}{s'}-\frac n{s}
=\frac{n'-n}{R}}.
\]
\excheck 公式的物像距离及曲率半径必须采用一致的有向正负约定；对 $n'=n$ 的界面，它退化为 $s'=s$ 而非产生折射光焦度。

\opticsproblem{30}{191}
图5.6：平行光从真空入射到折射率为 $n$ 的轴对称透明介质弯曲表面，欲使所有光线汇聚于顶点 O 后方轴向距离 $f$ 的 F 点。求界面形状 $x(r)$ 及可汇聚的最大入射束半径。
\exsolution
取 O 为 $(x,r)=(0,0)$，焦点 $F=(f,0)$。
与轴上入射光线相同光程的任一光线，先在真空中走到界面点 $(x,r)$，再在介质中到 F，满足
\[
x+n\sqrt{(f-x)^2+r^2}=nf.
\]
平方移项得到
$(n^2-1)x^2-2nf(n-1)x+n^2r^2=0$。
取 $x(0)=0$ 的根，
\[
\boxed{x(r)=\frac{nf}{n+1}
\left[1-\sqrt{1-\frac{(n+1)r^2}{(n-1)f^2}}\right]}.
\]
要求根号非负，因此
\[
\boxed{r_{\max}=f\sqrt{\frac{n-1}{n+1}}}.
\]
\excheck $r_{\max}$ 的长度量纲来自根号外的 $f$；书末答复的纯文本抽取对该根号不清晰，不宜误写为 $f(n-1)/(n+1)$。

\opticsproblem{31}{191}
点光源位于对称双凸玻璃厚透镜前表面
$20$cm 处；透镜厚 $d=5.0$cm，两面曲率半径的模均为
$R=5.0$cm。
求像位于透镜后表面之后多远。取普通玻璃
$n\simeq1.50$。
\exsolution
沿光传播方向取有向距离。
第一面从空气入玻璃，$R_1=+5$cm、物距 $s_1=-20$cm。
球面折射公式
$\frac{1.5}{s_1'}-\frac1{-20}
=\frac{0.5}{5}$（单位均取cm），
得到
$s_1'=30$cm。
第二面前的中间像位于其右方
$30-5=25$cm，故对第二面是虚物：
$s_2=+25$cm，
其曲率半径 $R_2=-5$cm，玻璃到空气：
\[
\frac1{s_2'}-\frac{1.5}{25}
=\frac{1-1.5}{-5}=0.10\,\mathrm{cm^{-1}},
\]
于是
$1/s_2'=0.16\,\mathrm{cm^{-1}}$，
\[
\boxed{s_2'=6.25\,\mathrm{cm}
\simeq6.3\,\mathrm{cm}}.
\]
\excheck 厚透镜必须分别处理前、后两个折射面，不能直接用薄透镜公式忽略5cm厚度。

\opticsproblem{32}{191--192}
图示厚 $d=9.0$cm 的平凸玻璃透镜，物体放在凸面前，
像恰形成于后方平面表面并以该面为屏。
凸面曲率半径 $R=2.5$cm，
物体亮度 $L=7700\,\mathrm{cd/m^2}$，
入瞳直径 $D=5.0$mm。
求：(a) 横向放大率；(b) 像照度。取玻璃折射率 $n=1.50$、无损耗。
\exsolution
设物体在空气中，进入玻璃后，像位于玻璃内距离前折射面 $d$ 处。
球面折射光学功率
$\Phi=(n-1)/R$，
后方像距 $s'=d$，
有向公式
$n/d-1/s=(n-1)/R$。
近轴横向放大率
$\beta=(n_{\rm in}s')/(n_{\rm out}s)
=d/(ns)$，
整理
\[
\boxed{\beta=1-\frac{d(n-1)}{nR}
=1-\frac{9(0.5)}{1.5(2.5)}
=-0.20}.
\]
以像方光束锥角和亮度守恒
$L_{\rm image}=n^2L$，
在小数值孔径近似下，
\[
\boxed{E\simeq\frac{\pi n^2D^2L}{4d^2}
\simeq42\,\mathrm{lx}}.
\]
\excheck 放大率负号表明实像倒立；亮度与面积比推导依赖近轴光束和忽略 Fresnel 损失。

\opticsproblem{33}{192}
求两种薄透镜的光焦度与物方、像方焦距：
(a) 玻璃透镜在空气中光焦度 $\Phi_0=-5.0$D，浸入
$n_0=1.7$ 的液体；
(b) 对称双凸玻璃薄透镜在空气中光焦度
$\Phi_0=+10$D，左侧为空气、右侧为水。
取玻璃折射率 $n\simeq1.50$、水折射率 $n_w\simeq1.33$。
\exsolution
(a) 相同面形的透镜光焦度与相对折射率差成正比：
\[
\Phi=\Phi_0\frac{n-n_0}{n-1}
=-5\frac{1.5-1.7}{1.5-1}=+2.0\,\mathrm D.
\]
均处于液体中，
$\boxed{f'=-f=n_0/\Phi=0.85\,\mathrm m}$。
(b) 两曲面光焦度之和
$\Phi=\frac{n-1}{R}+\frac{n_w-n}{-R}
=(2n-n_w-1)/R$；
空气中对称双凸 $\Phi_0=2(n-1)/R$，故
\[
\boxed{\Phi=\frac{\Phi_0(2n-n_w-1)}{2(n-1)}
\simeq6.7\,\mathrm D},
\]
$\boxed{f=-1/\Phi\simeq-15\,\mathrm{cm}}$，
$\boxed{f'=n_w/\Phi\simeq20\,\mathrm{cm}}$。
\excheck 双侧介质不同使物方与像方焦距绝对值不相同，不能照搬空气中 $f'=-f$。

\opticsproblem{34}{192}
按原书图5.7及图5.8作图：
(a) 会聚、发散薄透镜后的光线；
(b) 已知共轭点 P、P' 与主轴，求透镜位置和焦点；
(c) 已知某一条入射—出射光线对和主轴，求另一条入射光线的出射路径。
\exsolution
(a) 对会聚透镜，平行主轴入射光线出射后经过像方焦点 F'；
经过物方焦点 F 的入射光线出射后平行主轴；
通过光心 O 的近轴光线方向不变。
发散透镜应将“经过焦点”替换为“出射反向延长线经过虚焦点”。
(b) 连接给定共轭点 P、P'；这条不偏折的光心光线与主轴相交处即透镜光心，作与主轴垂直的透镜平面。
再用由 P 发出的平行主轴光线，结合 P' 连线与透镜平面交点求像方焦点 F'；倒作得 F。
(c) 将已知光线在透镜面上的交点及出射方向记录为主平面基准；
另一条入射光线可分解为穿过光心的不偏折光线与平行主轴光线，按共轭焦平面的同一点成像规则构造出射方向。
\excheck 本题是几何作图题而非数值题；给出了图5.7、5.8可操作的基准光线与焦平面构造规则，作图必须保持原图各射线位置关系。

\opticsproblem{35}{192}
焦距 $f=25$cm 的薄会聚透镜将物体像投影至距透镜 $l=5.0$m 的屏幕。屏幕向透镜方向移动
$\Delta l=18$cm。
问为使像重新清晰，物体应移动多少距离？
\exsolution
固定焦距，物距 $s$ 与像距 $l$ 满足
$1/f=1/s+1/l$，
即
$s(l)=fl/(l-f)$。
屏幕移至
$l'=l-\Delta l$，
物距新值
$s'=fl'/(l'-f)$，
因此精确位移为
\[
\Delta s=s'-s
=\frac{f^2\Delta l}{(l-f)(l-\Delta l-f)}.
\]
代 $f=0.25$m、$l=5.0$m、
$\Delta l=0.18$m，
\[
\boxed{\Delta s\simeq5.2\times10^{-4}\,\mathrm m
\approx0.5\,\mathrm{mm}}.
\]
\excheck 新像距变短后物距必须\emph{增加}，即物体远离透镜；近似 $\Delta s\simeq f^2\Delta l/(l-f)^2$。
